\documentclass[conference]{IEEEtran}
\IEEEoverridecommandlockouts
\usepackage{cite}
\usepackage{amsmath,amssymb,amsfonts}
\usepackage{algorithmic}
\usepackage{graphicx}
\usepackage{textcomp}
\usepackage{xcolor}
\usepackage{url}
\usepackage{hyperref}

\def\BibTeX{{\rm B\kern-.05em{\sc i\kern-.025em b}\kern-.08em
    T\kern-.1667em\lower.7ex\hbox{E}\kern-.125emX}}

\begin{document}


%===========================================
% Start your report from here.
%===========================================

\title{Natural Disaster Severity Prediction via\\
Gradient Boosting and Temporal Feature Engineering\\
{\footnotesize Data Mining Final Project, 2026 Spring}
}

\author{\IEEEauthorblockN{Hsin-Yu Chen}
\and
\IEEEauthorblockN{Wei-Hsin Hung}
\and
\IEEEauthorblockN{Sky Shih-Kai Hong}
}

\maketitle

\begin{abstract}
Group ID: 5. Members: Hsin-Yu Chen (314551141,
sychen.cs14@nycu.edu.tw), Wei-Hsin Hung (314553021,
velda.cs14@nycu.edu.tw), and Sky Shih-Kai Hong (314554024,
sky.cs14@nycu.edu.tw). Our code is available at GitHub:
\href{https://github.com/skyhong2002/disaster-severity-prediction}{\color{blue}https://github.com/skyhong2002/disaster-severity-prediction}.

Natural disaster severity prediction is a critical task for risk management
and early warning systems. In this work, we address the problem of forecasting
weekly drought severity scores (0--5 scale) for 2,248 geographic regions using
historical meteorological data. We employ a leakage-aware direct multi-step
forecasting approach using an ensemble of LightGBM, XGBoost, and CatBoost
models. The implemented system builds temporal features per horizon from
calendar variables, meteorological lags, rolling statistics, exponentially
weighted means, long-window drought proxies, climatology anomalies, and
91-day-gapped historical score context. By separating historical diagnostic
experiments from the current direct-horizon implementation, we avoid accidental
score leakage while preserving useful region-specific signals. The current best
reportable Kaggle public leaderboard score recorded by the repository is
\textbf{0.8124}. Separately, the selected already-submitted public-chase
artifact reaches \textbf{0.7915} and crosses Baseline 3 (\textbf{0.8056}); we
keep this final-selection artifact separate from the reportable method claim.
\end{abstract}

\section{Project Summary (\textcolor{red}{max 1 page})}

\subsection{Goal}
The goal of this project is to predict the drought severity score for each
geographic region over the \textbf{next five consecutive weeks} (35 days),
given 91 days of historical meteorological observations per region from the
test set. Scores range from 0 (no disaster) to 5 (extreme severity).

\subsection{High-Level Approach}
Our approach follows a supervised machine learning paradigm treating
the task as a \textbf{multi-step time series regression} problem. We extract
temporal features from the raw meteorological columns, including lag values,
rolling means/stds, exponentially weighted averages, calendar encodings,
climatology anomalies, and domain-specific drought indices. We also include
historical score features only after a 91-day gap, matching the blind test
window and preventing access to unavailable labels. These features feed into
direct horizon models using LightGBM, XGBoost, and CatBoost. To maximize
robustness, we ensemble their predictions. Final predictions are post-processed
by clipping values to the valid score interval [0, 5].

\subsection{Related Works}
Time series forecasting for environmental phenomena has been widely studied.
Gradient boosting methods, notably XGBoost~\cite{chen2016xgboost} and
LightGBM~\cite{ke2017lightgbm}, have demonstrated strong performance on
tabular time series data by capturing non-linear feature interactions without
requiring extensive feature scaling. For forecasting competitions and
meteorological applications, lag-window feature engineering is often effective
because recent observations and seasonal patterns usually contain strong
predictive signals~\cite{kaggle_m5}. Transformer-based models such as Temporal
Fusion Transformer~\cite{lim2021temporal} and PatchTST~\cite{nie2023patchtst}
are possible future extensions, but the current repository implementation
focuses on a fully reproducible gradient-boosted tree pipeline.

\section{Proposed Method (\textcolor{red}{1-3 pages})}

\subsection{Problem Formulation}
Let $\mathcal{R} = \{r_1, r_2, \ldots, r_K\}$ denote the set of geographic
regions. For each region $r_k$, we observe a time series of meteorological
features $\mathbf{X}_{r_k} \in \mathbb{R}^{T \times D}$ where $T$ is the
number of observed days and $D$ is the number of meteorological variables.
The target $y_{r_k, t}$ is the drought severity score at week $t$, available
only on a specific day per week (the rest are NaN).

Given the 91-day test window, the task is to predict:
\begin{equation}
  \hat{y}_{r_k, t+1}, \hat{y}_{r_k, t+2}, \ldots, \hat{y}_{r_k, t+5}
\end{equation}
\noindent minimizing the Mean Absolute Error:
\begin{equation}
  \text{MAE} = \frac{1}{K \cdot 5} \sum_{k=1}^{K} \sum_{w=1}^{5}
  \left| \hat{y}_{r_k, w} - y_{r_k, w} \right|
\end{equation}

\subsection{Data Preprocessing}

\subsubsection{Handling NaN Scores and Blind-Gap Features}
Scores are observed only once per week, so most daily rows have missing target
values. During training, only rows with real non-missing scores are used as
supervised labels. Earlier diagnostic experiments attempted to reconstruct
missing scores with a separate model, but that design created confusing
train-test behavior and is not the current implementation path. The current
pipeline keeps historical score context only when it is available before the
91-day blind gap. A lag of zero in the score-history feature group therefore
means ``last known score before the blind window,'' not the current or recent
test-period score.

\subsubsection{Feature Engineering}
For each meteorological column $d$, we construct the following feature groups:

\begin{itemize}
  \item \textbf{Meteorological lags}: 7, 14, 21, 28, 35, 42, and 49 days.
  \item \textbf{Rolling statistics}: shifted rolling mean and standard
        deviation over 7, 14, 21, 28, and 35 days.
  \item \textbf{EWM features}: exponentially weighted means with spans 3, 7,
        and 14 days.
  \item \textbf{Calendar features}: month, quarter, day of month, approximate
        day of year, week of year, and cyclical sine/cosine encodings.
  \item \textbf{Domain Features}: We introduced a Drought Index (temperature divided by precipitation) and a Dryness Index (daily temperature range multiplied by max temperature) to explicitly capture conditions favorable to drought.
  \item \textbf{Climatology anomalies}: region-month normal weather and
        standardized deviations for precipitation, temperature, humidity, and
        related signals.
  \item \textbf{Blind-gap score history}: score lags and rolling summaries
        shifted by at least 91 days to prevent leakage from the unlabeled test
        window.
  \item \textbf{Region statistics}: historical mean, standard deviation,
        median, maximum, and minimum score per region.
\end{itemize}

\subsection{Model Architecture}

\subsubsection{Direct Horizon Models}
We train one independent regressor per prediction horizon (week 1 through week 5),
a technique known as \textbf{Direct Multi-Step Forecasting}. Each horizon model
uses the full feature set and predicts exactly one target column,
\texttt{target\_w1} through \texttt{target\_w5}.

\subsubsection{Ensemble Strategy}
To capture diverse patterns and reduce variance, we train \textbf{LightGBM},
\textbf{XGBoost}, and \textbf{CatBoost} models. LightGBM uses leaf-wise growth,
XGBoost is configured with histogram-based depth-wise trees, and CatBoost uses
native categorical handling for features such as region and calendar buckets.
The current best reportable public submission uses a convex blend with weights
35\% LightGBM, 35\% XGBoost, and 30\% CatBoost.

\subsection{Inference and Submission Generation}
At inference time, recent training rows and the 91-day test rows are concatenated
and sorted by \texttt{region\_id} and \texttt{date}. We rebuild the feature table
using the feature options saved in the selected experiment run, select the last
row of the 91-day test window for each region, and apply the five saved horizon
models. The inference script can load LightGBM, XGBoost, or CatBoost model
pickles from a run directory. The predictions are then blended and outputted to
a timestamped CSV file under \texttt{submissions/}.

\subsection{Experiment Versioning}
Because the implementation is expected to evolve, each run is saved under a
timestamped directory \texttt{experiments/<run\_id>/}. The directory contains
\texttt{config.json} for hyperparameters and feature settings,
\texttt{metrics.json} for validation results, and
\texttt{submission\_metadata.json} for prediction statistics after inference.
Large model files and generated submission CSVs are stored inside the same run
directory but are ignored by git. This format allows LightGBM, XGBoost,
CatBoost, ensembles, and future neural time series models to be appended to the
report as comparable experiment rows. Some historical run identifiers still use
the string \texttt{two\_stage}; this is retained for artifact compatibility and
should not be interpreted as a description of the current training flow.

\section{Experiments (\textcolor{red}{2-4 pages})}

\subsubsection{Dataset and Experimental Setup}
The dataset contains meteorological measurements for multiple geographic
regions. The training set spans \textbf{5,480 days} per region with weekly
drought scores (0--5). The test set provides 91 days of observations per
region without labels; we must predict 5 future weekly scores.

Table~\ref{tab:data_stats} summarizes key dataset statistics.

\begin{table}[h]
\centering
\caption{Dataset Statistics}
\label{tab:data_stats}
\begin{tabular}{lcc}
\hline
\textbf{Property} & \textbf{Train} & \textbf{Test} \\
\hline
Days per region   & 5,480 & 91 \\
\# Regions        & 2,248 & 2,248 \\
\# Raw meteo features & 14 & 14 \\
\# Rows           & 12,319,040 & 204,568 \\
Score range       & 0--5  & N/A \\
Score NaN ratio   & 85.73\% & N/A \\
\hline
\end{tabular}
\end{table}

The implementation supports both a \textbf{Chronological Holdout Split} and a
\textbf{Rolling-Origin} validation mode. The holdout split is used for the
original LightGBM/XGBoost anchor runs, while rolling-origin validation is used
for the CatBoost robustness experiments. These validation scores should not be
compared as if they came from the same protocol; Kaggle public and private
leaderboard results remain the external check for generalization.

\subsubsection{Baseline Descriptions}
We compare against the official Kaggle baselines reported on the public
leaderboard. Baseline 1, Baseline 2, and Baseline 3 provide reference MAE
scores for evaluating whether our implemented pipeline improves over simple
competition-provided methods.

\subsubsection{Experimental Results}

Table~\ref{tab:experiment_log} records the main legal experiments and historical
diagnostics used to evaluate historical score features, validation behavior,
and model diversity.

\begin{table}[h]
\centering
\caption{Experiment Log \& Kaggle Leaderboard Evaluation}
\label{tab:experiment_log}
\begin{tabular}{lcc}
\hline
\textbf{Method Strategy} & \textbf{Local MAE} & \textbf{Public MAE} \\
\hline
Naive \texttt{ffill} Persistence (gap=0) & 0.2321 (Leaked) & 1.0815 \\
Strategy B (Pure Weather Direct) & 0.6770 & 0.8640 \\
Long-Term Weather (365 days) & 0.4888 & 0.8604 \\
LGB/XGB Anchor (50/50) & 0.6942 / 0.7150 & 0.8232 \\
CatBoost tail2737 (rolling-origin) & 0.2212 & component only \\
\textbf{LGB/XGB/CatBoost 35/35/30} & N/A & \textbf{0.8124} \\
Initial Submission (Magic Leak) & 0.2100 (Leaked) & 0.8094* \\
\hline
\multicolumn{3}{l}{\footnotesize *Discarded due to unintentional data leakage in feature engineering.}
\end{tabular}
\end{table}

Table~\ref{tab:leaderboard} reports the Kaggle public leaderboard snapshot
after our submissions. The public leaderboard is calculated with
approximately 40\% of the test data. For methodological reporting, the clean
Team 5 model lineage remains the 35\% LightGBM / 35\% XGBoost / 30\% CatBoost
ensemble with public MAE \textbf{0.8124}. The corresponding reportable artifact
is \path{submissions/ensemble_20260516_lgb_xgb_cat2737_35_35_30.csv}. For
Kaggle final-selection purposes, Team 5 also has an already-submitted
public-chase artifact with public MAE \textbf{0.7915}, which crosses Baseline 3.
This artifact is valid as a leaderboard selection candidate, but it is not used
as the primary reportable modeling claim.

\begin{table}[h]
\centering
\caption{Kaggle Public Leaderboard Snapshot}
\label{tab:leaderboard}
\begin{tabular}{lcc}
\hline
\textbf{Team / Baseline} & \textbf{Public MAE} & \textbf{Rank / Entries} \\
\hline
Team 3 & 0.7628 & 1 \\
Team 4 & 0.7817 & 2 \\
\textbf{Team 5 selected public-chase artifact} & \textbf{0.7915} & \textbf{3} \\
Team 20 & 0.7942 & 4 \\
Team 2 & 0.7955 & 5 \\
Team 1 & 0.8043 & 6 \\
Baseline 3 & 0.8056 & 7 \\
\textbf{Team 5 reportable LGB/XGB/CatBoost lineage} & \textbf{0.8124} & \textbf{--} \\
Baseline 2 & 0.8623 & 1 \\
Baseline 1 & 0.9117 & 1 \\
\hline
\multicolumn{3}{l}{\footnotesize Public-chase artifact is selected for leaderboard status, not method reporting.}
\end{tabular}
\end{table}

The selected public-chase row corresponds to
\path{submissions/baseline3_private_hedge_v3_cat35_lower_w1_more_horizon_0p225_0p375_0p55_0p75_1.csv}
(Kaggle ref \texttt{53074980}, SHA-12 \texttt{003351676087}). It was selected
after public-feedback probing as a horizon-specific private-hedge frontier
artifact, so it should not be described as a new training algorithm. It tied
ref \texttt{53075001} at public MAE \textbf{0.7915}, but was preferred because
it is closer to the clean reportable anchor. More anchor-tilted submitted hedge
alternatives include same-day ref \texttt{53075022} and cross-day refs
\texttt{53038040}, \texttt{53038036}, and \texttt{53038033}. A teammate
LightGBM v3 enhanced candidate was also screened and submitted as ref
\texttt{53074655}, but its public MAE \textbf{1.0685} made it a negative
public readout rather than a reportable contribution. After the 2026-05-27
quota was exhausted, we prepared an unsubmitted v4 week-1 frontier queue for
the next reset; these files have no leaderboard score yet and are not method
claims. The reportable method lineage remains the 35/35/30
LGB/XGB/CatBoost artifact with SHA-12 \texttt{bee6f618828d}.

The generated submission contains one row per region and five predicted scores.
Table~\ref{tab:pred_stats} summarizes the prediction distribution for the current
best reportable public submission. The mean forecast gradually decreases from week 1
to week 5, while all predictions remain inside the valid score range after
clipping.

\begin{table}[h]
\centering
\caption{Prediction Statistics for Current Best Legal Submission}
\label{tab:pred_stats}
\begin{tabular}{lcccc}
\hline
\textbf{Column} & \textbf{Mean} & \textbf{Std.} & \textbf{Min} & \textbf{Max} \\
\hline
\texttt{pred\_week1} & 0.955 & 0.819 & 0.000 & 3.470 \\
\texttt{pred\_week2} & 0.927 & 0.818 & 0.000 & 3.376 \\
\texttt{pred\_week3} & 0.880 & 0.824 & 0.000 & 3.457 \\
\texttt{pred\_week4} & 0.831 & 0.804 & 0.000 & 3.428 \\
\texttt{pred\_week5} & 0.817 & 0.827 & 0.000 & 3.583 \\
\hline
\end{tabular}
\end{table}

\subsubsection{Ablation Study: The Kaggle Paradox \& Data Leakage}
To systematically analyze the impact of historical score dependencies and identify potential data leakage, we conducted an extensive ablation study.

\textbf{1. The ``Magic Leak'' Phenomenon (v0):} Our initial submission achieved a surprisingly strong Public MAE of 0.8094. However, investigation revealed this was caused by a train-test discrepancy where the model inadvertently relied on recent scores during training, while the testing phase relied on a fixed forward-filled (\texttt{ffill}) score. This created an accidental hybrid persistence model. To maintain scientific rigor, we discarded this model despite its high score.

\textbf{2. The Naive Persistence Failure (1.0815):} We explicitly reproduced the pure \texttt{ffill} behavior by setting the score gap to 0. While this yielded an artificially low Local MAE of 0.2321 (due to looking at the current day's true score), it failed spectacularly on Kaggle with an MAE of 1.0815. This proves that drought is highly dynamic and cannot be modeled by simply freezing the last known state for 91 days.

\textbf{3. The Kaggle Paradox \& Long-Term Weather (0.8604):} We initially hypothesized that removing the noisy historical scores entirely (Strategy B) would improve generalization. We expanded this by providing the model with up to 365 days of accumulated weather data, successfully dropping the Local MAE to an impressive 0.4888. However, the Kaggle Public MAE remained at 0.8604. This paradox demonstrates that historical scores, even when noisy, contain crucial unobserved confounders (e.g., local water infrastructure, soil characteristics) that pure meteorological features cannot capture.

\textbf{4. The CatBoost Diversity Gain (0.8124):} After establishing the legal
LGB/XGB anchor, we added a time-aware CatBoost direct-horizon model using native
categorical handling and rolling-origin validation. A 35/35/30 blend of
LightGBM, XGBoost, and CatBoost improved the legal public MAE from 0.8232 to
0.8124. Because the CatBoost validation protocol differs from the original
holdout protocol, we treat the public improvement as promising but still require
private leaderboard evidence before over-tuning its weight.

% ========================================================
\section{Conclusion}
In this work, we implemented a robust time series prediction framework for forecasting drought severity levels. By using chronological and rolling-origin validation modes and by enforcing 91-day-gapped score-history features, we reduced leakage risk while retaining useful regional severity context. Our ablation study revealed a compelling ``Kaggle Paradox'': pure long-term meteorological models perform strongly locally but fail to generalize to the leaderboard because they lack unobserved geographical and infrastructural confounders implicitly encoded in historical drought scores. We demonstrated that naive target persistence fails completely (MAE 1.0815), while the current direct-horizon ensemble of LightGBM, XGBoost, and CatBoost achieves a competitive legal Public MAE of 0.8124. For final leaderboard selection, an already-submitted public-chase artifact reaches public MAE 0.7915 and crosses Baseline 3, but we treat it as a selected Kaggle artifact rather than evidence of a new reportable training method.

% ========================================================
\bibliographystyle{IEEEtran}
\bibliography{reference}

\end{document}
